% =====================================================================
%  Reference documentation for interface/ibpy.py
%
%  Build (minted needs shell-escape + Pygments on PATH):
%     # make sure `pygmentize` is reachable, e.g. the project venv:
%     export PATH="$PWD/../../.venv2/bin:$PATH"
%     pdflatex -shell-escape ibpy_documentation.tex
%     pdflatex -shell-escape ibpy_documentation.tex   # 2nd pass for ToC
% =====================================================================
\documentclass[10pt,a4paper]{article}

\usepackage[T1]{fontenc}
\usepackage[utf8]{inputenc}
\usepackage[margin=2.2cm]{geometry}
\usepackage{lmodern}
\usepackage{microtype}
\usepackage{xcolor}
\usepackage{minted}
\usepackage{enumitem}
\usepackage{seqsplit}
\usepackage{titlesec}
\usepackage{fancyhdr}
\usepackage{hyperref}

\hypersetup{
    colorlinks=true, linkcolor=blue!55!black, urlcolor=blue!55!black,
    pdftitle={interface/ibpy.py -- function reference},
}

% ---- minted setup ----
\setminted[python]{
    fontsize=\small,
    breaklines=true,
    breakanywhere=false,
    frame=single,
    framesep=5pt,
    rulecolor=\color{gray!40},
    bgcolor=gray!7,
    tabsize=4,
    autogobble,
}
\newminted{python}{}                       % environment: pythoncode
% inline code: robust verbatim that also works inside \item[...] labels,
% and that may line-break (long signatures otherwise overflow the margin).
% detokenize is pre-expanded into \ibpytmp so seqsplit sees raw characters.
\newcommand{\py}[1]{%
  \begingroup
    \edef\ibpytmp{\detokenize{#1}}%
    \texttt{\expandafter\seqsplit\expandafter{\ibpytmp}}%
  \endgroup}

\definecolor{secblue}{RGB}{40,82,140}
\titleformat{\section}{\large\bfseries\color{secblue}}{\thesection}{0.6em}{}
\titleformat{\subsection}{\normalsize\bfseries}{\thesubsection}{0.6em}{}
\titlespacing*{\section}{0pt}{14pt}{6pt}

\setlist[description]{leftmargin=1.4em,style=nextline,itemsep=2pt,parsep=1pt}
\newcommand{\fn}[1]{\py{#1}}               % function-name shortcut

\pagestyle{fancy}
\fancyhf{}
\lhead{\small\textsc{interface/ibpy.py}}
\rhead{\small function reference}
\cfoot{\thepage}
\renewcommand{\headrulewidth}{0.3pt}

% =====================================================================
\begin{document}

\begin{titlepage}
\centering
\vspace*{3cm}
{\Huge\bfseries The \texttt{ibpy} Interface}\\[10pt]
{\LARGE A reference to \texttt{interface/ibpy.py}}\\[24pt]
{\large The single boundary between this animation package and Blender's
\texttt{bpy} API}\\[40pt]
\begin{minipage}{0.82\textwidth}\small
\itshape
\texttt{ibpy.py} encapsulates almost all direct calls into Blender. Every
scene in the \py{video_*} folders talks to Blender \emph{through} this
module, never to \texttt{bpy} directly. The practical payoff: when the
Blender API changes between versions, the fix lives in one file; and the
module doubles as a searchable catalogue of ``how do I do X in this
package''. This document mirrors the section ordering that the source file
itself now uses --- search the source for \texttt{=== <SECTION NAME>} to
jump straight to the corresponding code.
\end{minipage}
\vfill
{\small Generated for the reorganised \texttt{ibpy.py}
(375 functions, 25 sections).}
\end{titlepage}

\tableofcontents
\newpage

% =====================================================================
\section{How to read this document}
% =====================================================================

The module is a flat collection of free functions; you import it and call
them with a Blender object (conventionally a \py{bob} --- ``Blender
object'') plus parameters:

\begin{pythoncode}
from interface import ibpy
# or:  from interface.ibpy import set_camera_location, camera_zoom
\end{pythoncode}

A few conventions recur throughout:

\begin{description}
\item[\texttt{bob}] the target object. Most functions accept either a real
  \py{bpy} object or one of this package's wrapper objects and call
  \fn{get_obj} internally to normalise it.
\item[\py{begin_time} / \py{transition_time}] animation timing in
  \emph{seconds}. They are converted to frames with the module-level
  \py{FRAME_RATE}. Many ``\py{change_*}'' functions return
  \py{begin_time + transition_time}, so you can chain them:
  \py{t0 = ibpy.change_default_value(..., begin_time=t0)}.
\item[\py{begin_frame} / \py{frame_duration}] the older,
  frame-based variants of the same idea (integer frames).
\item[\texttt{slot}] index into an object's material slots (default
  \py{0}).
\end{description}

The sections below follow the file's own ordering. Within each section the
most-used functions are documented with runnable usage snippets taken from
the \py{video_*} scenes; the remaining helpers are summarised so the
section is complete.

% =====================================================================
\section{Context, scene, frames \& file I/O}
% =====================================================================

Entry points to Blender's global state and to saving/loading.

\begin{description}
\item[\fn{get_context()}, \fn{get_scene()}, \fn{get_layout()}] thin
  accessors for \py{bpy.context}, the active scene, and the ``Layout''
  screen. Used internally far more than in scene code.
\item[\fn{set_frame(frame)}, \fn{get_frame()}, \fn{start_frame()},
  \fn{end_frame()}] move/query the playhead and the scene's frame range.
\item[\fn{close(x, y, precision=EPS)}] floating-point ``almost equal''
  test used by geometry helpers.
\item[\fn{save(filename)}, \fn{load(filepath)}, \fn{load_file(filename)},
  \fn{append(filename)}] persistence. \fn{append} pulls assets from an
  external \texttt{.blend}.
\item[\fn{set_cursor_location(position)}, \fn{cursor_to_origin()}] drive
  the 3D cursor (which several ``set origin'' operations rely on).
\item[\fn{register_class(cls)}] register a custom \py{bpy} class.
\end{description}

\begin{pythoncode}
# advance the timeline and persist the result
ibpy.set_frame(0)
ibpy.save("/tmp/scene_v1.blend")
\end{pythoncode}

% =====================================================================
\section{Object lookup, creation, copying \& deletion}
% =====================================================================

\begin{description}
\item[\fn{get_obj(bob)}] the workhorse: returns the underlying \py{bpy}
  object for a wrapper, a name, or a \py{bpy} object (idempotent). Almost
  every other function calls it first.
\item[\fn{get_obj_from_name(name)}, \fn{get_objects_from_name(name)}] look
  up by (sub)string of the object name.
\item[\fn{rename(bob, name)}] rename an object.
\item[\fn{create(mesh, name, location)}] wrap a mesh datablock into a new
  scene object at a location.
\item[\fn{add_empty(**kwargs)}] add an Empty (handy as a camera target or
  parent pivot).
\item[\fn{copy(bob)}, \fn{duplicate(bob)}] shallow copy / full duplicate.
\item[\fn{delete(obj)}] remove an object from the scene.
\item[\fn{get_child_with_name(bob, child_name, starts_with=False)}] find a
  specific child in a hierarchy.
\end{description}

\begin{pythoncode}
target = ibpy.add_empty(location=[0, 0, -6], name="CameraTarget")
clone  = ibpy.duplicate(some_object)
ibpy.rename(clone, "Clone.001")
\end{pythoncode}

% =====================================================================
\section{Mesh primitives \& mesh creation}
% =====================================================================

Convenience wrappers around Blender's primitive operators, plus the
all-important \fn{create_mesh} for arbitrary geometry.

\begin{description}
\item[\fn{create_mesh(vertices, edges=[], faces=[], name='mesh',
  orient_faces=False)}] build a mesh datablock from raw
  vertex/edge/face lists. The most general entry point; everything else is
  a shortcut.
\item[\fn{add_sphere}, \fn{add_cube}, \fn{add_cone}, \fn{add_cylinder},
  \fn{add_torus}, \fn{add_circle}, \fn{add_plane}] the usual primitives;
  all take \py{location}, \py{scale}, and \py{smooth} where meaningful.
\item[\fn{add_reference_image(name, **kwargs)}] add an image-empty for
  modelling reference.
\item[\fn{create_vertex_group(bob, vertex_indices, name)},
  \fn{add_vertex_group(bob, indices, name, weight_function)}] define
  weighted vertex groups (used by modifiers and deformations).
\end{description}

\begin{pythoncode}
# build a polyhedron from explicit data and hand it to a wrapper object
mesh = ibpy.create_mesh(vertices, faces=faces)
kwargs_red = {"mesh": mesh, "color": "trunc_icosidodeca"}

# a smooth-shaded reference sphere
ball = ibpy.add_sphere(radius=1, location=(0, 0, 0), resolution=5,
                       smooth=True)
\end{pythoncode}

% =====================================================================
\section{Mesh editing, conversion \& analysis}
% =====================================================================

\begin{description}
\item[\fn{get_vertex_locations(bob)}, \fn{get_vertices(bob)}] read back
  vertex coordinates / vertices.
\item[\fn{smooth_mesh(bob)}, \fn{shade_smooth(auto=True)}] smoothing.
\item[\fn{analyse_mesh(bob)}] print a structural report (vertex / edge /
  face counts) --- a debugging aid.
\item[\fn{apply(bob, location=True, rotation=True, scale=True)}] apply
  transforms into the mesh data (bakes them into vertex coordinates).
  \fn{apply_scale(bob)} is the scale-only shortcut.
\item[\fn{convert_to_mesh(obj, apply_transform=False)}] convert a curve /
  text / other object into a mesh.
\item[\fn{separate(bob, type='SELECTED')}] split selected geometry into a
  new object.
\item[\fn{select_faces(faces, normal)},
  \fn{stretch_uv_map_to_faces(bob, normal)}] face selection and UV
  stretching (used for image-on-geometry materials).
\item[\fn{add_attribute(bob, name, attribute, type='FLOAT',
  domain='POINT')}] add a custom mesh attribute (read later by geometry
  nodes).
\end{description}

\begin{pythoncode}
ibpy.apply(bob, location=False, rotation=True, scale=True)
verts = ibpy.get_vertex_locations(bob)     # list of Vectors
ibpy.add_attribute(bob, name="charge", attribute=charges, type="FLOAT")
\end{pythoncode}

% =====================================================================
\section{Selection \& edit/object mode}
% =====================================================================

A complete set of (de)selection helpers plus mode switching. They wrap the
fiddly \py{bpy.ops} selection state machine.

\begin{description}
\item[\fn{set_edit_mode()}, \fn{set_object_mode()}] toggle interaction
  mode.
\item[\fn{set_active(obj)}, \fn{get_active()}] manage the active object.
\item[\fn{select(obj, deselect_others=True)}, \fn{un_select(obj)},
  \fn{deselect_all()}] basic selection.
\item[\fn{select_recursively(obj)}, \fn{select_all_deep(bobs)},
  \fn{deep_select(obj)}] select an object together with its descendants.
\item[\fn{set_select(obj, value=True)},
  \fn{set_several_select(objects, value=True)}] bulk set selection flags
  without operators.
\end{description}

\begin{pythoncode}
ibpy.deselect_all()
ibpy.select(bob)                 # also makes it active, deselecting others
ibpy.set_edit_mode()
# ... operate ...
ibpy.set_object_mode()
\end{pythoncode}

% =====================================================================
\section{Collections \& linking}
% =====================================================================

\begin{description}
\item[\fn{make_new_collection(name, hide_render=False,
  hide_viewport=False)}] create a collection.
\item[\fn{get_collection(name)}, \fn{get_collection_name(obj)},
  \fn{remove_collection(collection)}] manage collections.
\item[\fn{link(obj, collection=None, recursively=True)},
  \fn{recursive_link(obj, collection)}, \fn{un_link(obj, collection)}]
  link/unlink objects into a collection (objects must be linked to appear
  in the scene). \fn{is_linked(obj)} tests membership.
\item[\fn{to_collection(collection)},
  \fn{collection_to_string(collection)}] utilities used by the importer.
\end{description}

\begin{pythoncode}
coll = ibpy.make_new_collection("Particles", hide_render=False)
ibpy.link(particle_obj, collection="Particles", recursively=True)
\end{pythoncode}

% =====================================================================
\section{Visibility (hide / unhide) \& fading}
% =====================================================================

Two flavours of ``make it disappear'': \emph{hard} hide (toggles
\py{hide_render}) and \emph{soft} fade (animates material alpha). The
\py{*_rec} / \py{recursive_*} variants walk the child hierarchy.

\begin{description}
\item[\fn{hide(bob, begin_time=0)}, \fn{unhide(bob, begin_time=0)}]
  keyframe render visibility at a time.
\item[\fn{hide_rec}, \fn{unhide_rec}, \fn{hide_frm_rec}, \fn{unhide_frm_rec}]
  recursive / frame-based versions.
\item[\fn{is_hidden(obj)}, \fn{set_hide(bob, hide, frame)}] query / set raw
  flags.
\item[\fn{fade_in(b_obj, frame, frame_duration)},
  \fn{fade_out(b_obj, frame, frame_duration, alpha=0, children=True)}]
  animate alpha up/down; recurse into children by default.
  \fn{fade_out_quickly} and \fn{recursive_fade_out} are variants.
\item[\fn{set_shadow(value)}, \fn{set_shadow_of_object(b_object, shadow)}]
  toggle shadow casting.
\end{description}

\begin{pythoncode}
ibpy.fade_in(label, frame=0, frame_duration=30)
# later: hide from render after 5 s
ibpy.hide(label, begin_time=5)
\end{pythoncode}

% =====================================================================
\section{Lights \& light probes}
% =====================================================================

\begin{description}
\item[\fn{add_point_light(**kwargs)}, \fn{add_area_light(**kwargs)},
  \fn{add_spot_light(energy, **kwargs)}] add lights; kwargs include
  \py{location}, \py{radius}, \py{scale}, \py{energy}.
\item[\fn{set_sun_light(location, energy=2)}, \fn{remove_sun_light()}] the
  scene sun.
\item[\fn{add_light_probe(**kwargs)}] add a reflection/irradiance probe.
\item[\fn{switch_on(bob, ...)}, \fn{switch_off(bob, ...)},
  \fn{change_power(bob, from_value, to_value, ...)}] animate a light's
  power; \fn{get_energy_at_frame(bob, frame)} reads it back.
\item[\fn{light_path_settings(total=12, diffuse=4, glossy=4,
  transmission=12, volume=0, transparency=8)}] global Cycles light-path
  bounce limits --- a common quality/perf knob.
\end{description}

\begin{pythoncode}
ibpy.add_point_light(location=[4, -4, 6], energy=2000, radius=0.5)
ibpy.light_path_settings(total=12, transmission=12, transparency=8)
ibpy.change_power(lamp, from_value=0, to_value=5000,
                  begin_frame=0, frame_duration=30)
\end{pythoncode}

% =====================================================================
\section{Camera setup \& animation}
% =====================================================================

The most frequently used section. A scene typically: adds/places a camera,
points it at a target Empty, sets the lens, then animates zoom/move.

\begin{description}
\item[\fn{add_camera(location, rotation)}, \fn{get_camera()}] create / get
  the camera.
\item[\fn{set_camera_location(location, frame=0)},
  \fn{set_camera_rotation(rotation)}] place it.
\item[\fn{set_camera_lens(lens=50, clip_end=1000)}] focal length and far
  clip.
\item[\fn{set_camera_view_to(target, rotation_euler=[0,0,0],
  targetZ=False, up_axis='UP_Y')}] aim the camera at a target with a
  Track-To constraint --- the standard way to frame a subject.
\item[\fn{camera_zoom(lens, begin_time, transition_time)}] animate the
  focal length (dolly-zoom feel).
\item[\fn{camera_move(shift, begin_time, transition_time)}] translate the
  camera over time.
\item[\fn{camera_rotate_to(rotation_euler, ...)}] animate orientation.
\item[\fn{set_camera_dof(focus_object, aperture_fstop=0.1)}] depth of
  field.
\item[\fn{set_camera_follow(target)}, \fn{set_camera_copy_location(target)},
  \fn{camera_follow(target, initial_value, final_value, ...)},
  \fn{camera_change_track_influence(...)},
  \fn{camera_change_follow_influence(...)}] constraint-based camera rigs
  (follow a path, copy a location, blend influence in/out).
\item[\fn{camera_alignment_euler(bob, camera_location)},
  \fn{camera_alignment_quaternion(bob, camera_location, default)}] compute
  the rotation that makes an object face the camera (billboarding).
\end{description}

\begin{pythoncode}
# canonical opening shot (from video_4d_geometry)
camera_empty = ibpy.add_empty(location=[0, 0, -6])
ibpy.set_camera_view_to(camera_empty)
ibpy.set_camera_location(location=Vector([-30, 55, 55]))
ibpy.set_camera_lens(lens=50)

# animate a slow zoom-in over 2 s
ibpy.camera_zoom(lens=30, begin_time=0.1, transition_time=2)
# then track the action for the rest of the shot
ibpy.camera_move(shift=[0, 0, 4], begin_time=2, transition_time=6)
\end{pythoncode}

% =====================================================================
\section{Constraints, tracking, follow paths \& parenting}
% =====================================================================

The general constraint machinery that the camera section builds on, plus
object parenting.

\begin{description}
\item[\fn{add_constraint(b_object, type, name=None, **kwargs)}] add any
  Blender constraint by type string.
\item[\fn{set_track(bob, target)}, \fn{set_copy_location(bob, target)},
  \fn{set_follow(bob, target, **kwargs)}] add Track-To / Copy-Location /
  Follow-Path constraints.
\item[\fn{follow(bob, target, initial_value, final_value, begin_time,
  transition_time)}] animate movement along a follow-path target.
\item[\fn{set_track_influence}, \fn{set_follow_influence},
  \fn{change_follow_influence}] blend constraint influence over time.
\item[\fn{set_parent(b_obj, parent)}, \fn{get_parent(bob)},
  \fn{get_oldest_parent(bob)}, \fn{clear_parent(bob)}] object parenting.
\item[\fn{set_vertex_parent(parent, child, vertex_index)}] parent a child
  to a single vertex (used for the ``chain of arrows'' constructions).
\end{description}

\begin{pythoncode}
ibpy.set_parent(child_obj, parent_obj)        # rigid hierarchy
ibpy.set_track(arrow, target=tip_empty)       # always point at the tip
\end{pythoncode}

% =====================================================================
\section{Transform animation (move / rotate / scale / grow / shrink)}
% =====================================================================

Keyframed transforms. Frame-based functions take \py{begin_frame,
frame_duration}; the \fn{shrink}/\fn{grow} pair has both a time-based and a
scale-based overload (the later definition wins, see the note at the end).

\begin{description}
\item[\fn{move(b_obj, direction, begin_frame, frame_duration)},
  \fn{move_to(b_obj, target, begin_frame, frame_duration,
  global_system=False)}, \fn{move_fast_from_to(b_obj, start, end, ...)}]
  translate by a delta / to an absolute point.
\item[\fn{rotate_by(b_obj, rotation_euler)},
  \fn{rotate_to(b_obj, begin_frame, frame_duration, pivot,
  interpolation)}] rotation (instant / animated about a pivot).
\item[\fn{grow(b_obj, scale, begin_frame, frame_duration,
  initial_scale=0, modus='from_center')},
  \fn{shrink(b_obj, scale, ...)}] scale-in / scale-out animation.
\item[\fn{grow_from(b_obj, pivot, ...)}, \fn{shrink_from(b_obj, pivot,
  ...)}, \fn{rescale(b_obj, re_scale, ...)}] scale about an explicit
  pivot.
\item[\fn{set_location(b_obj, location)}, \fn{set_origin(bob,
  type='ORIGIN_GEOMETRY')}, \fn{set_pivot(obj, origin)},
  \fn{transform_into_world(b_obj, location)}] non-animated placement and
  origin control.
\end{description}

\begin{pythoncode}
# scale an object up from nothing over 0.5 s (15 frames @ 30 fps)
ibpy.grow(bob, scale=1, begin_frame=0, frame_duration=15,
          initial_scale=0, modus="from_center")

ibpy.set_origin(bob, type="ORIGIN_GEOMETRY")   # center the pivot
\end{pythoncode}

% =====================================================================
\section{Materials: assignment, color, alpha \& emission}
% =====================================================================

The largest section (37 functions). Covers attaching materials to objects
and animating their scalar/colour properties.

\begin{description}
\item[\fn{add_material(bob, material)}, \fn{set_material(bob, material,
  slot=0)}, \fn{append_materials(bobj, materials)}] attach materials to
  slots. \fn{get_material(material, **kwargs)} resolves a material by name
  or spec; \fn{get_materials(bob)}, \fn{get_material_of(bob, slot)} read
  them back.
\item[\fn{assign_material_to_faces(bob, material_index, normal)},
  \fn{set_color_to_faces(bob, face2color_function)},
  \fn{create_color_map_for_mesh(bob, colors, name)},
  \fn{set_vertex_colors(bob, colors)}] per-face / per-vertex colouring.
\item[\fn{change_alpha(b_obj, frame, frame_duration, ...)},
  \fn{change_alpha_of_material(mat, from_value, to_value, begin_time,
  transition_time)}, \fn{set_alpha_for_material(material, alpha)},
  \fn{set_alpha_and_keyframe(obj, value, frame, slot)}] animate
  transparency. \fn{get_alpha_at_current_keyframe(obj, frame, slot)} reads
  it.
\item[\fn{set_emission}, \fn{set_emission_color}, \fn{change_emission},
  \fn{change_emission_to}, \fn{change_emission_by_name},
  \fn{change_emission_of_material}] drive emission strength / colour.
\item[\fn{set_metallic_for_material}, \fn{set_roughness_for_material},
  \fn{set_transmission_for_material},
  \fn{set_emission_strength_for_material}] one-shot PBR property setters.
\item[\fn{change_brightness}, \fn{change_color}, \fn{mix_color},
  \fn{shader_value}, \fn{adjust_mixer}, \fn{set_mixer}] colour-mix and
  value-node animation. \fn{get_alpha_mixer},
  \fn{change_material_properties}, \fn{get_new_material} are construction
  helpers.
\end{description}

\begin{pythoncode}
ibpy.set_material(bob, my_material, slot=0)

# fade a material to fully transparent over 1 s
ibpy.change_alpha_of_material(mat, from_value=1, to_value=0,
                              begin_time=t0, transition_time=1)

# pulse the emission of slot 0
ibpy.change_emission(bob, from_value=0, to_value=5,
                     slot=0, begin_frame=0, frame_duration=30)
\end{pythoncode}

% =====================================================================
\section{Shader / geometry node access \& default-value drivers}
% =====================================================================

Find nodes inside a material or node group, and animate their socket
default values --- the backbone of procedural animation in this package.

\begin{description}
\item[\fn{get_node_tree(name, type='GeometryNodeTree')}] create/fetch a
  node tree.
\item[\fn{get_node_with_label(bob, label)},
  \fn{get_node_from_tree(tree, label)},
  \fn{get_node_from_shader(shader, label)},
  \fn{get_all_nodes_from_shader(shader, label)},
  \fn{get_shader_node_from_material(material, label)},
  \fn{get_nodes_of_material(bob, name_part, slot)},
  \fn{get_node_of_material(material, label)}] locate nodes by label.
\item[\fn{change_default_value(slot, from_value, to_value, begin_time,
  transition_time)}] the single most useful driver: keyframe a socket's
  \py{default_value} from one number to another. Returns the end time for
  chaining.
\item[\fn{set_default_value}, \fn{change_default_boolean},
  \fn{set_default_boolean}, \fn{change_default_integer},
  \fn{change_default_vector}, \fn{change_default_rotation},
  \fn{change_default_quaternion}, \fn{change_default_material},
  \fn{change_value}] type-specific variants for each socket kind.
\item[\fn{change_shader_value(b_object, node, input, initial_value,
  final_value, frame, frame_duration)}, \fn{change_mixer(mixer, ...)}]
  animate a specific node input / a mix factor.
\item[\fn{get_default_value_at_frame(socket, frm)},
  \fn{get_value_at_frame(function, frm)},
  \fn{get_parameter_at_frame(function, frm)}] sample animated values.
\item[\fn{add_item_to_switch(switch_node, idx, socket_or_value, tree)},
  \fn{remove_unlinked_nodes(node_tree)}] node-tree editing utilities.
\end{description}

\begin{pythoncode}
# chain several socket animations, threading the time t0 through:
t0 = 0
t0 = 0.5 + ibpy.change_default_value(sliders[-1], from_value=0,
                                     to_value=1, begin_time=t0,
                                     transition_time=1)
t0 = ibpy.change_default_value(ra, from_value=1, to_value=0,
                               begin_time=t0, transition_time=1)
\end{pythoncode}

% =====================================================================
\section{Procedural material \& shader builders, textures}
% =====================================================================

Functions that \emph{construct} entire shader node graphs --- the
domain-specific materials used across the videos (Mandelbrot colouring,
hue maps, magnet fields, gradients, image/movie textures).

\begin{description}
\item[\fn{customize_material(material, **kwargs)}] the central material
  customiser; pass \py{override_material} to replace special materials.
\item[\fn{make_mandelbrot_material}, \fn{make_hue_material},
  \fn{make_iteration_material}, \fn{make_magnet_material},
  \fn{make_gradient_material}, \fn{make_dashed_material},
  \fn{make_alpha_frame}] ready-made procedural materials.
\item[\fn{make_coarse_grained_group}, \fn{make_phase2color_group},
  \fn{make_vector_interpolation}, \fn{make_mandel_iteration}] reusable node
  sub-groups consumed by the materials above.
\item[\fn{make_image__stretched_over_geometry(src, v_min, v_max)},
  \fn{add_image_texture(bob, image, ...)},
  \fn{setup_material_for_alpha_masking(bob)},
  \fn{create_image_mixing_find_previous_mixer(material, image_path)}]
  image-texture pipeline (incl.\ alpha masking).
\item[\fn{set_movie_to_material(bob, src, duration, ...)},
  \fn{set_movie_start(bob, begin_frame)}] map a video clip onto a surface.
\item[\fn{create_color_mixing(material, color1, color2)},
  \fn{create_color_mixing_find_previous_color(material, color)},
  \fn{create_shader_from_script(material, osl_script, ...)},
  \fn{create_voronoi_shader(material, colors, ...)},
  \fn{create_shader_from_function(material, hue_functions, ...)}] colour
  mixers and OSL/Voronoi/RPN shader generators.
\end{description}

\begin{pythoncode}
mat = ibpy.make_hue_material()
ibpy.set_material(plane, mat)
ibpy.add_image_texture(plane, image="diagram.png",
                       begin_time=0, transition_time=1)
\end{pythoncode}

% =====================================================================
\section{Node-group \& RPN function builders}
% =====================================================================

Generic machinery that turns a mathematical formula in \emph{reverse Polish
notation} into a Shader/Geometry node group. These power
\fn{make_function}-style construction throughout the package.

\begin{description}
\item[\fn{create_group_from_scalar_function(nodes, functions,
  parameters=[], ...)}, \fn{create_group_from_vector_function(nodes,
  functions, ...)}] build a node group computing scalar / vector RPN
  formulae.
\item[\fn{create_shader_group_from_function(nodes, function, parameters,
  inputType, outputType, name)}] shader-tree variant.
\item[\fn{create_iterator_group(nodes, functions, parameters, iterations,
  name)}] nest a function group \py{iterations} times (fractal / iterative
  maps).
\item[\fn{if_node(nodes, bool_function, parameters, ...)}] a mix node
  switching between two inputs based on a boolean RPN expression.
\item[\fn{make_new_socket(tree, name, io='INPUT',
  type='NodeSocketFloat')}] add a typed socket to a group interface (used
  everywhere groups are built).
\item[\fn{build_group_component}, \fn{build_scalar_group_component},
  \fn{build_group}, \fn{create_part}, \fn{build_recursively}] the recursive
  RPN-to-nodes compiler internals. \fn{build_group} is deprecated in favour
  of \fn{build_group_component}.
\end{description}

\begin{pythoncode}
# add a vector input socket to a freshly created group tree
ibpy.make_new_socket(tree, name="position", io="INPUT",
                     type="NodeSocketVector")
\end{pythoncode}

% =====================================================================
\section{Geometry-nodes modifier access \& sockets}
% =====================================================================

Reach into a geometry-nodes modifier to read its nodes and push values into
its exposed sockets.

\begin{description}
\item[\fn{get_geometry_nodes_modifier(bob)}] return the geometry-nodes
  modifier on an object. \emph{(Defined twice; the second definition ---
  the documented ``extract modifier from bob'' one --- wins.)}
\item[\fn{get_geometry_node_from_modifier(modifier, label)}] fetch an
  inner node by label, e.g.\ to drive it later.
\item[\fn{get_socket_names_from_modifier(modifier)}] list the sockets that
  receive custom values during setup. \emph{(Also defined twice; last
  wins.)}
\item[\fn{set_socket_data_for_geometry_node_modifier(bob, socket_data)}]
  bulk-assign exposed-socket values.
\item[\fn{get_material_from_modifier(modifier, label)}] pull a material
  generated inside the modifier.
\end{description}

\begin{pythoncode}
mod        = ibpy.get_geometry_nodes_modifier(bob)
csv_node   = ibpy.get_geometry_node_from_modifier(mod, "Import CSV")
counter    = ibpy.get_geometry_node_from_modifier(mod, "Counter")
ibpy.change_default_value(counter.inputs["Value"], from_value=0,
                          to_value=100, begin_time=0, transition_time=4)
\end{pythoncode}

% =====================================================================
\section{Volume scatter \& absorption}
% =====================================================================

\begin{description}
\item[\fn{set_volume_scatter(bob, value, begin_time=0)},
  \fn{change_volume_scatter(bob, final_value, begin_time,
  transition_time)}] animate volumetric scattering on an object.
\item[\fn{set_volume_absorption(bob, value, begin_time=0)},
  \fn{change_volume_absorption(bob, final_value, ...)}] volumetric
  absorption.
\item[\fn{set_volume_scatter_of_material(material, value)},
  \fn{set_volume_absorption_of_material(material, value)}] material-level
  setters.
\item[\fn{set_volume_scatter_background(density=1)}] world-volume fog.
\end{description}

\begin{pythoncode}
ibpy.change_volume_scatter(cloud, final_value=0.4,
                           begin_time=t0, transition_time=2)
\end{pythoncode}

% =====================================================================
\section{World, background, HDRI, render \& compositor}
% =====================================================================

Render-engine configuration, the world/HDRI environment, and post-render
compositor effects. \fn{set_render_engine} and \fn{set_hdri_background} are
among the most-called functions in the whole package.

\begin{description}
\item[\fn{set_render_engine(engine='CYCLES', transparent=False,
  motion_blur=False, denoising=False, resolution_percentage=100,
  frame_start=0, taa_render_samples=1024, ...)}] one call to configure
  the render pipeline.
\item[\fn{set_hdri_background(filename, ext='exr', simple=False,
  transparent=False, strength=0.2, rotation_euler=None, reflections=True,
  ...)}] load an HDRI as the world background.
\item[\fn{set_hdri_strength(strength, begin_time, transition_time)},
  \fn{change_hdri_strength(start, end, ...)},
  \fn{hdri_background_rotate(rotation_euler, ...)}] animate the HDRI.
\item[\fn{get_background()}, \fn{get_sky()}, \fn{animate_sky_background()},
  \fn{dim_background(value, ...)},
  \fn{set_volume_scatter_background(density)}] other world controls.
\item[\fn{create_composition(denoising=None)}, \fn{empty_blender_view3d()},
  \fn{set_taa_render_samples(taa_render_samples=64, begin_frame=0)},
  \fn{set_exposure(value)}] compositor / viewport / exposure.
\item[\fn{animate_glare(glare, **kwargs)},
  \fn{change_glow_threshold(...)}, \fn{change_glow_size(...)},
  \fn{get_image(src)}] glare/glow compositor effects.
\end{description}

\begin{pythoncode}
# canonical render + world setup (from video_4d_geometry)
ibpy.set_hdri_background("kloofendal_misty_morning_puresky_4k", 'exr',
                         simple=True, transparent=True,
                         rotation_euler=Vector())
t0 = ibpy.set_hdri_strength(1, begin_time=0, transition_time=0)
ibpy.set_render_engine(denoising=False, transparent=True, frame_start=1,
                       resolution_percentage=100, engine="CYCLES",
                       taa_render_samples=64, motion_blur=True)
\end{pythoncode}

% =====================================================================
\section{Curves, splines \& bevel}
% =====================================================================

Bezier-curve construction and the bevel/extrude controls that turn a curve
into a renderable tube, plus ``grow along the curve'' animation.

\begin{description}
\item[\fn{import_curve(path)}, \fn{new_curve_object(name, curve)},
  \fn{get_new_curve(name, num_points, data, **kwargs)}] curve creation /
  import.
\item[\fn{add_bezier_spline(splines, num_points=2, data, cyclic=True)},
  \fn{add_points_to_spline(spline, total_points, closed_loop)},
  \fn{set_bezier_point_of_curve(curve, index, position, handle_left,
  handle_right)}, \fn{reset_bezier_point(...)}] build / edit splines and
  control points.
\item[\fn{get_splines(obj)}, \fn{get_splines_of_curve(b_obj)},
  \fn{get_bezier_points(spline)},
  \fn{get_curve_for_object(obj)}, \fn{get_all_curves()},
  \fn{merge_splines(b_objects)}, \fn{separate_pieces(data)},
  \fn{get_spline_bounding_box(b_obj)}] inspect / combine curve data.
\item[\fn{set_bevel(bob, depth, caps=True, res=4)},
  \fn{change_bevel(bob, old_depth, new_depth, ...)},
  \fn{set_extrude(bob, extrude)},
  \fn{set_bevel_factor_mapping(bob, ...)}] give a curve thickness.
\item[\fn{grow_curve(bob, appear_frame, duration, ...)},
  \fn{shrink_curve(bob, appear_frame, duration, inverted=False)},
  \fn{set_curve_range(bob, factor, ...)},
  \fn{set_curve_full_range(bob, factor1, factor2, ...)},
  \fn{set_use_path(bob, is_used)}] animate a curve drawing itself in /
  out (via bevel-factor start/end).
\end{description}

\begin{pythoncode}
ibpy.set_bevel(curve_obj, depth=0.02, caps=True, res=6)
# draw the curve on, left to right, over 1.5 s
ibpy.grow_curve(curve_obj, appear_frame=0, duration=45)
\end{pythoncode}

% =====================================================================
\section{Shape keys \& morphing}
% =====================================================================

Animate between mesh shapes via shape keys --- the basis of the ``morph
one object into another'' transitions.

\begin{description}
\item[\fn{add_shape_key(bobj, name, previous=None, following=None,
  relative=True)}] add a shape key.
\item[\fn{create_shape_key_from_transformation(bob, old_sk, index,
  transformation)},
  \fn{create_shape_key_from_index_transformation(bob, old_sk, index,
  transformation)}] derive a new shape by transforming an existing one.
\item[\fn{keyframe_shape_key(bob, sk_name, frame)},
  \fn{set_shape_key_to_value(bob, sk_name, value, frame)},
  \fn{set_shape_key_eval_time(bob, time, frame)}] keyframe shape-key
  values / evaluation time.
\item[\fn{morph_to_next_shape}, \fn{morph_to_next_shape2},
  \fn{morph_to_next_shape_with_pause}, \fn{morph_to_previous_shape},
  \fn{morph_to_first_shape}, \fn{set_to_first_shape},
  \fn{set_to_last_shape}, \fn{zero_all_shape_keys}] high-level morph
  sequencing (the \py{*2} variant zeroes all but the current key).
\end{description}

\begin{pythoncode}
ibpy.add_shape_key(bob, "target")
ibpy.morph_to_next_shape(bob, current_shape_index=0,
                         appear_frame=0, frame_duration=30)
\end{pythoncode}

% =====================================================================
\section{Modifiers}
% =====================================================================

Add, configure, animate and apply Blender modifiers (incl.\ geometry
nodes).

\begin{description}
\item[\fn{add_mesh_modifier(bob, **kwargs)}] the generic, preferred
  entry point (dispatches on a \py{type} kwarg). \fn{replace_mesh_modifier}
  swaps one out.
\item[\fn{add_modifier(b_obj, type, name=None)},
  \fn{add_modifier_recursive(obj, type, name)}] low-level add by type.
\item[\fn{add_sub_division_surface_modifier(b_obj, level=2, ...)},
  \fn{add_solidify_modifier(b_obj, thickness=0.1)},
  \fn{add_bevel_modifier(b_obj, width=0.1)}] common specific modifiers.
\item[\fn{change_solidifier_thickness(bob, from_value, to_value, ...)},
  \fn{change_modifier_attribute(bob, mod_name, attr_name, start, end,
  ...)}] animate modifier parameters.
\item[\fn{apply_modifier(bob, type=None)}, \fn{apply_modifiers(bob)}] bake
  modifiers into the mesh.
\end{description}

\begin{pythoncode}
ibpy.add_solidify_modifier(shell, thickness=0.05)
ibpy.change_modifier_attribute(bob, "Solidify", "thickness",
                               start=0, end=0.1,
                               begin_frame=0, frame_duration=30)
ibpy.apply_modifiers(bob)
\end{pythoncode}

% =====================================================================
\section{Keyframes, F-curves \& animation interpolation}
% =====================================================================

The low-level keyframe layer that all the \py{change_*} functions sit on,
plus tools to make animations linear and to debug F-curves.

\begin{description}
\item[\fn{insert_keyframe(bob, data_path, frame)}] insert a keyframe on
  any data path.
\item[\fn{set_bevel_factor_and_keyframe(data, value, frame)},
  \fn{set_bevel_factor_start_and_end_keyframe(data, start, end, frame)},
  \fn{set_bevel_factor_start_keyframe(data, start, frame)}] curve-draw
  keyframing primitives used by \fn{grow_curve}.
\item[\fn{set_linear_fcurves(bob)}, \fn{set_linear_action(bob,
  data_path_part)}, \fn{set_linear_action_full(bob)},
  \fn{set_linear_action_modifier(bob)}, \fn{set_bezier_action_modifier(bob)},
  \fn{set_linear_fcurves_for_nodes(node)},
  \fn{make_animations_linear(thing, data_paths, extrapolate=False)}]
  force linear (constant-speed) interpolation on selected channels.
\item[\fn{clear_animation_data(bob)}] wipe all animation from an object.
\item[\fn{value_of_action_fcurve_at_frame(action, fcurve, frm)},
  \fn{print_actions()}, \fn{print_fcurves_of_action(action)}] inspection /
  debugging helpers (``first used in \py{video_apollonian}'').
\end{description}

\begin{pythoncode}
ibpy.insert_keyframe(bob, "location", frame=0)
ibpy.make_animations_linear(bob, data_paths=["location"])
\end{pythoncode}

% =====================================================================
\section{Rigid body, forces \& physics simulation}
% =====================================================================

\begin{description}
\item[\fn{make_rigid_body(bob, dynamic=True, kinematic=False, friction=0.1,
  bounciness=0.5, ...)}, \fn{make_rigid_bodies(bob_list, ...)}] turn one
  or many objects into rigid bodies.
\item[\fn{add_rigid_body(bob, **kwargs)},
  \fn{key_frame_rigid_body_properties(bob, type='dynamic', value=True,
  begin_time=0)}] lower-level rigid-body control and keyframing the
  dynamic/kinematic switch.
\item[\fn{add_force(**kwargs)}, \fn{add_wind(**kwargs)},
  \fn{add_turbulence(**kwargs)}, \fn{decay_force(bob, initial_strength,
  final_strength=0, ...)}] force fields.
\item[\fn{set_simulation(begin_time=0, transition_time=250/FRAME_RATE)}]
  run/scope the physics simulation over a time window.
\end{description}

\begin{pythoncode}
for child in pieces:
    ibpy.make_rigid_body(child, type="ACTIVE",
                         collision_shape="MESH", friction=0.5)
ibpy.add_force(location=[0, 0, 0], strength=-50)   # gravity well
ibpy.set_simulation(begin_time=0, transition_time=4)
\end{pythoncode}

% =====================================================================
\section{Geometry queries: bounding box, centers, world location}
% =====================================================================

Read-only spatial queries: where is an object, how big is it, where is a
particular feature. Heavily used to position labels and arrows relative to
existing geometry.

\begin{description}
\item[\fn{get_bounding_box(bob)}, \fn{get_bounding_box_for_letter(letter)},
  \fn{analyse_bound_box(bob)}, \fn{get_dimension(bob, direction)}] extents
  and size.
\item[\fn{find_center_of_leftest_vertices(obj)} and its
  \texttt{rightest / lowest / highest / closest / furthest} siblings] the
  geometric centre of an object's extremal face/edge/vertex --- ideal
  anchor points for annotations.
\item[\fn{get_world_location(b_obj)}, \fn{get_world_location2(b_obj)},
  \fn{get_world_matrix(bob)}, \fn{get_location(b_obj)},
  \fn{get_scale(b_obj)}, \fn{get_rotation(b_obj)}] current transform.
\item[\fn{get_world_location_at_time(b_obj, time)},
  \fn{get_location_at_frame(b_obj, frame)},
  \fn{get_scale_at_frame(b_obj, frame)},
  \fn{get_rotation_at_frame(b_obj, frame)},
  \fn{get_rotation_quaternion_at_frame(b_obj, frame)},
  \fn{get_offset_factor_at_frame(b_obj, target, frame)},
  \fn{get_input_value_at_frame(b_obj, frame)}] sample the transform /
  values at a specific time --- essential when one animation must start
  where another leaves off.
\end{description}

\begin{pythoncode}
top = ibpy.find_center_of_highest_vertices(bob)   # anchor a label here
where = ibpy.get_world_location_at_time(mover, time=2.0)
\end{pythoncode}

% =====================================================================
\section{Miscellaneous utilities}
% =====================================================================

\begin{description}
\item[\fn{to_vector(z)}] coerce a list/tuple into a \py{mathutils.Vector}
  (\py{to_vector([1,2,3]) -> Vector((1.0, 2.0, 3.0))}).
\item[\fn{vectorize(list)}] map a nested list to Vectors.
\item[\fn{diff(a, b)}] element-wise difference helper.
\end{description}

\begin{pythoncode}
v = ibpy.to_vector([1, 2, 3])     # Vector((1.0, 2.0, 3.0))
\end{pythoncode}

% =====================================================================
\section{Notes on duplicate definitions}
% =====================================================================

Three names are defined twice in the source; Python keeps the
\emph{last} definition, and the reorganisation preserves that ``last wins''
ordering so behaviour is unchanged:

\begin{itemize}[leftmargin=1.4em]
\item \fn{get_geometry_nodes_modifier} --- the second (documented
  ``extract modifier from bob'') definition is the live one.
\item \fn{get_socket_names_from_modifier} --- second definition wins.
\item \fn{shrink} --- the scale-based overload (with a \py{scale}
  argument) wins over the earlier time-based one.
\end{itemize}

If you intend to call the \emph{earlier} behaviour, rename one of the
definitions; calling the name as-is always reaches the later one.

\vfill
\begin{center}\small\itshape
Generated from the reorganised \texttt{interface/ibpy.py}. Section order and
titles match the in-file \texttt{===} banners; search the source for a
banner to jump to the code.
\end{center}

\end{document}
